If $G$ is a subgroup of $S_m$ and $H$ a subgroup of $S_n$, we shall write $G{\wr}H$ for the subgroup of $S_m{\wr}S_n$
consisting of all elements $(\sigma;\alpha_1,\alpha_2,\ldots,\alpha_n)$ for $\alpha_1,\alpha_2,\ldots,\alpha_n \in G$ and $\sigma \in H$.
We shall make frequent use of such groups where $G$ and $H$ are each either the full symmetric group or a Young subgroup thereof, and we
shall often restrict or induce modules between such groups, for example
$\left.X\right\uparrow^{k(S_m{\wr}S_n)}_{k(S_m{\wr}S_\gamma)}$ and $\left.Y\right\downarrow^{k({S_m{\wr}S_n})}_{k({S_m{\wr}S_\gamma})}$
where $\gamma$ is some composition of $n$. As with the symmetric group, we shall de-clutter such expressions where possible by suppressing
the field and replacing the full symbols for subgroups of $S_m$ and $S_n$ with the subscript used to identify them,
so for example the above would be abbreviated to
$\left.X\right\uparrow^{m{\wr}n}_{m{\wr}\gamma}$ and $\left.Y\right\downarrow^{m{\wr}n}_{m{\wr}\gamma}$.

We now extend the notion of a Young subgroup to encompass multicompositions.
Let $\ul{\gamma} = (\gamma^1,\ldots,\gamma^t)$ be a $t$-multicomposition of $n$ ($t$ some non-negative integer) and let $\hat\gamma$ be
the composition of $n$ obtained by concatenating the compositions $\gamma^1,\ldots,\gamma^t$ in that order (so $\hat\gamma$ consists of
the parts of $\gamma^1$, followed by the parts of $\gamma^2$, and so on). We define the \textit{Young subgroup} of $S_n$ associated to
$\ul{\gamma}$ to be the Young subgroup $S_{\hat\gamma}$ associated to $\hat\gamma$, and we write $S_{\ul{\gamma}}$ for this subgroup. Thus we
have a canonical isomorphism $S_{\ul{\gamma}} \cong S_{\gamma^1} \times S_{\gamma^2} \times \cdots \times S_{\gamma^t}$.
Further, we note that $S_{\ul{\gamma}}$ is a subgroup of $S_{|\ul{\gamma}|}$.

We now recall several standard methods for constructing modules for wreath products, as described in
\cite[Section 4.3]{JAMKER} and \cite[Section 3]{CHTAN}. Recall that we are using \textit{right} modules.

Firstly, let $G$ be a subgroup of $S_m$, and
let $X$ be a $kG$-module. We define $X^{\widetilde\boxtimes n}$ to be the $k(G{\wr}S_n)$-module obtained by
equipping the $k$-vector space $X^{\otimes n}$ (that is, the tensor product over $k$ of $n$ copies of $X$) with the action given by the
formula
\[
	(x_1\otimes\cdots\otimes x_n)(\sigma;\alpha_1,\ldots,\alpha_n) = (x_{(1)\sigma^{-1}}\alpha_1)\otimes\cdots
	\otimes(x_{(n)\sigma^{-1}}\alpha_n)
\]
for $x_1,\ldots,x_n \in X$, $\alpha_1,\ldots,\alpha_n \in G$, $\sigma \in S_n$.
More generally, let $X_1,\ldots,X_t$ be $kG$-modules, and $\gamma = (\gamma_1,\ldots,\gamma_t)$ a composition of $n$ of length $t$.
We form a $k(G{\wr}S_{\gamma})$-module by equipping the $k$-vector space
$\bigl(X^{\otimes \gamma_1}_1\bigr)\otimes\bigl(X^{\otimes \gamma_2}_2\bigr)\otimes\cdots\otimes\bigl(X^{\otimes \gamma_t}_t\bigr)$
with the action given by the formula
\[
	(x_1\otimes\cdots\otimes x_n)(\sigma;\alpha_1,\ldots,\alpha_n) = (x_{(1)\sigma^{-1}}\alpha_1)\otimes\cdots
	\otimes(x_{(n)\sigma^{-1}}\alpha_n)
\]
where each $x_i$ lies in the appropriate $X_j$, $\alpha_1,\ldots,\alpha_n \in G$, and $\sigma \in S_{\gamma}$. We denote this module by
$\bigl(X_1,\ldots,X_t\bigr)^{\widetilde\boxtimes \gamma}$, and we note that $X^{\widetilde\boxtimes n}$ is the
special case of this construction where $\gamma$ has an $n$ in one place and all the other parts are 0.

Now let $G$ be a subgroup of $S_m$, $H$ a subgroup of $S_n$, and $Y$ a $kH$-module. It is easy to check that we may make $Y$ into a
$k(G{\wr}H)$-module via the formula
\begin{equation}\label{y_inf_act:eq}
y(\sigma;\alpha_1,\ldots,\alpha_n) = y\sigma
\end{equation}
for $y \in Y$, $\alpha_1,\ldots,\alpha_n \in G$, and $\sigma \in H$. This module may be understood by noting that $G{\wr}H$ is the
semidirect product of the normal subgroup consisting of all elements $(e;\alpha_1,\ldots,\alpha_n)$ for $\alpha_1,\ldots,\alpha_n \in G$
with the subgroup consisting of all elements $(\sigma;e,\ldots,e)$ for $\sigma \in H$. This latter subgroup is canonically isomorphic to
$H$, and hence we see that the module obtained from $Y$ via \eqref{y_inf_act:eq} is the \textit{inflation} of $Y$ from $H$ to $G{\wr}H$
with respect to the semidirect product
structure. Hence we shall denote this module by $\mathrm{Inf}^{G{\wr}H}_H Y$.
Now let $H$ be a subgroup of $S_n$, $G$ be a subgroup of $S_m$, $Y$ be a $kH$-module, and further let $Z$ be a $k(G{\wr}H)$-module. Then
we define a $k(G{\wr}H)$-module $Z{\oslash}Y$ as follows: the underlying $k$-vector space is $Z{\otimes}Y$, and
the action is given by the formula
\[(z \otimes y)(\sigma;\alpha_1,\ldots,\alpha_n) = (z(\sigma;\alpha_1,\ldots,\alpha_n))\otimes(y\sigma)\]
for $z \in Z$, $y \in Y$, $\alpha_1,\ldots,\alpha_n \in G$, $\sigma \in H$.
Thus we see that we have an equality of $k(G{\wr}H)$-modules $Z{\oslash}Y = Z \otimes \mathrm{Inf}^{G{\wr}H}_H Y$
where the module on the right-hand side is the internal tensor product of the $k(G{\wr}H)$-modules $Z$ and $\mathrm{Inf}^{G{\wr}H}_H Y$.
Since taking the (internal) tensor product of group modules and inflating group modules are both exact functors, it follows that the
operation $- \oslash -$ ``preserves submodule series'' in both places, in the sense that if $Z$ has a submodule series with quotients
$X_1,\ldots,X_t$ in order from top to bottom, then $Z \oslash Y$ will have a submodule series with quotients $X_1 \oslash Y, \ldots,
X_t \oslash Y$ in order from top to bottom, with a symmetrical result for a submodule series for $Y$.

We can combine the above constructions as follows: if $G$ is a subgroup of $S_m$, $X_1,\ldots,X_t$ are $kG$-modules and $Y$ is a
$kS_{\gamma}$-module for $\gamma$ a composition of $n$, then we obtain a $k(G{\wr}S_{\gamma})$-module
$\bigl(X_1,\ldots,X_t\bigr)^{\widetilde\boxtimes\gamma}{\oslash}Y$
with underlying vector space
$\bigl(X^{\otimes \gamma_1}_1\bigr)\otimes\bigl(X^{\otimes \gamma_2}_2\bigr)\otimes\cdots\otimes\bigl(X^{\otimes \gamma_t}_t\bigr){\otimes}Y$
and action given by the formula
\begin{multline}\label{xs_twid_osl_y_act:eq}
	(x_1{\otimes}\cdots{\otimes}x_n{\otimes}y)(\sigma;\alpha_1,\ldots,\alpha_n) =
	(x_{(1)\sigma^{-1}}\alpha_1)\otimes\cdots\otimes(x_{(n)\sigma^{-1}}\alpha_n){\otimes}(y\sigma)
\end{multline}
for $x_i \in X$, $\alpha_i \in G$, $y \in Y$, $\sigma \in S_{\gamma}$.

We now recall an elementary construction for producing $kS_{\gamma}$-modules $Y$ for use in the above constructions. Indeed, for each
$i \in \{1,\ldots,t\}$, let $Y_i$ be a right $kS_{\gamma_i}$-module. Now recall that we have a canonical identification of the group
$S_{\gamma}$ with the direct product $S_{\gamma_1}\times S_{\gamma_2}\times\cdots\times S_{\gamma_t}$ of groups. Thus any module for
$k(S_{\gamma_1}{\times}S_{\gamma_2}{\times}\cdots{\times}S_{\gamma_t})$ may be regarded as a $kS_{\gamma}$-module in a canonical way, and vice
versa. In particular, if $Y_i$ is a $kS_{\gamma_i}$-module for each $i$, then the external tensor product
$Y_1 \boxtimes Y_2 \boxtimes \cdots \boxtimes Y_t$, which is a
$k\left(S_{\gamma_1}\times S_{\gamma_2}\times\cdots\times S_{\gamma_t}\right)$-module, may be regarded as a $kS_{\gamma}$-module.

Now if $G$ is a subgroup of $S_m$ and $\gamma$ is a composition of $n$, then we have an obvious isomorphism between $G{\wr}S_{\gamma}$ and 
$\left(G{\wr}S_{\gamma_1}\right)\times\left(G{\wr}S_{\gamma_2}\right)\times\cdots\times\left(G{\wr}S_{\gamma_t}\right)$, and hence we have a
canonical identification of algebras between $k(G{\wr}S_{\gamma})$ and
$k\left(G{\wr}S_{\gamma_1}\right)\otimes k\left(G{\wr}S_{\gamma_2}\right)\otimes\cdots \otimes k\left(G{\wr}S_{\gamma_t}\right)$.
With this identification, it is now easy to see that we have an isomorphism of modules
\begin{multline}\label{mod_wprd_etp_vrs:eq}
\bigl(X_1,\ldots,X_t\bigr)^{\widetilde\boxtimes\gamma} {\oslash}\bigl(Y_1 \boxtimes Y_2 \boxtimes \cdots\boxtimes Y_t\bigr)
\cong\\
	\bigl(X_1^{\widetilde\boxtimes \gamma_1}\oslash Y_1\bigr)\boxtimes
	\bigl(X_2^{\widetilde\boxtimes \gamma_2}\oslash Y_2\bigr)\boxtimes\cdots
	\boxtimes\bigl(X_t^{\widetilde\boxtimes \gamma_t}\oslash Y_t\bigr)
\end{multline}
(this isomorphism was given in \cite[Lemma 3.2 (1)]{CHTAN}).

\begin{prop}\label{res_comm_tilde_boxtimes:prop}
Let $G_1 \subseteq G_2$ be subgroups of $S_m$ and $X$ a $kG_2$-module. Then we have an isomorphism of $k(G_1 \wr S_n)$-modules
\[
\Bigl[X^{\widetilde\boxtimes n}\Bigr]\Bigr\downarrow^{G_2 \wr S_n}_{G_1 \wr S_n} \:\cong\:
\Bigl[X\bigr\downarrow^{G_2}_{G_1}\Bigl]^{\widetilde\boxtimes n}.
\]
\end{prop}
\begin{proof}
This is immediate from the definition of $(-)^{\widetilde\boxtimes n}$.
\end{proof}

\begin{prop}\label{wreath_ind_res:prop}\cite[Lemma 3.2]{CHTAN}
Let $G$ be a subgroup of $S_m$.
Let $\alpha = (\alpha_1,\ldots,\alpha_t)$ be a composition of $n$ and let $V$ be a $k(G{\wr}S_n)$-module, $W$ a
$k(G{\wr}S_\alpha)$-module, $X$ a $kS_n$-module, and $Y$ a $kS_\alpha$-module. Then we have module isomorphisms
\begin{enumerate}
\item $\left.\bigl[V \oslash X\bigr]\right\downarrow^{G{\wr}n}_{G{\wr}\alpha} \cong
\bigl(\left.V\right\downarrow^{G{\wr}n}_{G{\wr}\alpha}\bigr) \oslash \left(\left.X\right\downarrow^{n}_{\alpha}\right)$
\item $V \oslash \left(\left.Y\right\uparrow^{n}_{\alpha}\right)
\cong \left.\bigl[\bigl(\left.V\right\downarrow^{G{\wr}n}_{G{\wr}\alpha}\bigr) \oslash Y\bigr]\right\uparrow^{G{\wr}n}_{G{\wr}\alpha}$
\item $\bigl(\left.W\right\uparrow^{G{\wr}n}_{G{\wr}\alpha}\bigr) \oslash X
\cong \left.\bigl[W \oslash \left(\left.X\right\downarrow^{n}_{\alpha}\right)\bigr]\right\uparrow^{G{\wr}n}_{G{\wr}\alpha}$
\end{enumerate}
where the symbols $n$ and $\alpha$ represent the subgroups $S_n$ and $S_\alpha$ of $S_n$, respectively.
\end{prop}

\section{Wreath product Specht modules}
We now define analogues for the wreath product $S_m \wr S_n$ of the Specht modules of the symmetric group using the above constructions.
As mentioned in the introduction, although these constructions are well-known, the author is not aware that these modules have
previously been considered as analogues of the symmetric group Specht modules.

Firstly we define some useful notation. If $Y_1,\ldots,Y_s$ are $kS_m$-modules and $\ul{\eta} = (\eta^1,\ldots,\eta^s)$ is an
$s$-component multipartition of $n$, then we define the $k(S_m{\wr}S_n)$-module $S^{\ul{\eta}}(Y_1,\ldots,Y_s)$ by setting
\[
S^{\ul{\eta}}(Y_1,\ldots,Y_s)\;=\; \left.\left[\bigl(Y_1,\ldots,Y_s\bigr)^{\widetilde\boxtimes |\ul{\eta}|}{\oslash}
\bigl(S^{\eta^1}\boxtimes\cdots\boxtimes S^{\eta^s}\bigr)\right]\right\uparrow^{m{\wr}n}_{m{\wr}|\ul{\eta}|}.
\]

We take $r$ to be the number of distinct partitions of $m$, and we enumerate them in the lexicographic order as follows
\[ (m) = \mu^1 > \mu^2 > \ldots > \mu^r = (1^m). \]
Then for $\ul{\nu}=(\nu^1,\ldots,\nu^r)$ an $r$-multipartition of $n$, we define a $k(S_m{\wr}S_n)$-module
\[
	S^{\ul{\nu}} =S^{\ul{\nu}}(S^{\mu^1},\ldots,S^{\mu^r})
\]
and we call $S^{\ul{\nu}}$ the the \textit{Specht module} for $S_m{\wr}S_n$ associated to $\ul{\nu}$. For later convenience, we also
define a $k(S_m \wr S_{|\ul{\nu}|})$-module
\[
	T^{\ul{\nu}} = \bigl(S^{\mu^1},\ldots,S^{\mu^r}\bigr)^{\widetilde\boxtimes |\ul{\nu}|}{\oslash}
	\bigl(S^{\nu^1}\boxtimes\cdots\boxtimes S^{\nu^r}\bigr)
\]
so that $S^{\ul{\nu}} = T^{\ul{\nu}}\bigr\uparrow^{m{\wr}n}_{m{\wr}|\ul{\nu}|}$.
As mentioned above, the use of the name ``Specht module'' here is justified by the fact that these modules have the same relationship
with the simple modules of $k(S_m \wr S_n)$ as the Specht modules of $kS_n$ have with the simple of $kS_n$. This may be demonstrated by
noting that the wreath product Specht modules occur as the cell modules of a cellular structure on $k(S_m \wr S_n)$, in the sense of
Graham and Lehrer \cite{GLCA}. For details, see \cite{GCSWPA} and \cite[Section 5.5]{PHDTHES}.

We now consider how filtrations of the $kS_m$-modules $Y_1,\ldots,Y_s$ induce filtrations of the $k(S_m \wr S_n)$-module
$S^{\ul{\eta}}(Y_1,\ldots,Y_s)$. This question was answered by Chuang and Tan
in \cite{CHTAN}, and the results we now present are taken from there. However, we shall present these results in a very slightly modified
form, using the notion of a \textit{multipartition matrix}, which is simply a matrix whose entries are multipartitions. We shall typically
denote the multipartition matrix whose $(i,j)$\tss{th} entry is the multipartition $\ul{\epsilon}^{ij}$ as $[\ul{\epsilon}]$. Thus a
multipartition matrix is simply a matrix whose entries are tuples of tuples of integers. Now let $s$ and $t$ be positive integers,
let $\alpha,\beta$ be compositions of the same integer $n$ and with lengths $s$ and $t$ respectively, and let $L$ be an $s{\times}t$
matrix with non-negative integer entries. We define $\mathrm{Mat}_{\ul{\Lambda}}(L;\alpha{\times}\beta)$ to be the
set of all $s{\times}t$ multipartition matrices $[\ul{\epsilon}]$ such that:
\begin{enumerate}
\item for each $i=1,\ldots,s$, the sum of all of the integers occurring in the $i$\tss{th} row of $[\ul{\epsilon}]$ is equal to the
$i$\tss{th} part of $\alpha$;
\item for each $j=1,\ldots,t$, the sum of all of the integers occurring in the $j$\tss{th} column of $[\ul{\epsilon}]$ is equal to the
$j$\tss{th} part of $\beta$;
\item the length of the $(i,j)$\tss{th} entry of $[\ul{\epsilon}]$ is equal to the $(i,j)$\tss{th} entry of $L$.
\end{enumerate}
Note that we do allow entries of $L$ to be zero, meaning that the corresponding entry of $[\ul{\epsilon}]$ is the empty multipartition
$()$. Note also that we allow multipartitions to contain the empty partition $()$ as an entry, which means, for example, that we can have
a multipartition of length 1 but size 0 in $[\ul{\epsilon}]$, namely the multipartition $(())$ whose sole entry is an empty partition.

For $[\ul{\epsilon}] \in \mathrm{Mat}_{\ul{\Lambda}}(L;\alpha{\times}\beta)$, we define $R_i[\ul{\epsilon}]$ to be the multipartition
obtained by concatenating the multipartitions on the $i$\tss{th} row of $[\ul{\epsilon}]$, and we similarly define $C_j[\ul{\epsilon}]$
to be the multipartition obtained by concatenating the multipartitions on the $j$\tss{th} column of $[\ul{\epsilon}]$ (ordered from top
to bottom).

From \cite{CHTAN} we have the following result. Now \cite{CHTAN} formally works only with the semisimple case, but the arguments given
there work in the general case, if one speaks of filtrations rather than direct sum decompositions and makes use of the  modular
Littlewood--Richardson result provided by \cite{JAMPEEL}. Further, \cite{CHTAN} formally makes the assumption that the modules
$X_1,\ldots,X_t$ are pairwise non-isomorphic, but this is not in fact needed for the proof if we speak in terms of filtrations
\textit{admitting} multiplicities (as explained in Section~\ref{filt_ser:subsec}) to allow for non-uniqueness. For a very detailed proof
of the result in essentially the form below (taking account of the aforementioned minor discrepancies between our situation and that of
\cite{CHTAN}), see Section 6.4 of the author's PhD thesis, \cite{PHDTHES}.

\begin{prop}\label{S_eta_Y_filt_S_nu_X:prop}\textup{(}\cite[Lemma 4.4, (1)]{CHTAN}; see also \cite[Proposition 6.4.1]{PHDTHES}\textup{)}
Let $Y_1,\ldots,Y_s$ and $X_1,\ldots,X_t$ be $kS_m$-modules such that for each $i=1,\ldots,s$, $Y_i$ has a filtration by the modules
$X_1,\ldots,X_t$ admitting the multiplicities $(a^i_j)_j$ (where we allow $a^i_j=0$). Let $\ul{\eta}$ be an $s$-component
multipartition of $n$. Then $S^{\ul{\eta}}(Y_1,\ldots,Y_s)$ has a filtration by the modules $S^{\ul{\nu}}(X_1,\ldots,X_t)$ for $\ul{\nu}$ a
$t$-multipartition of $n$, admitting the multiplicities $\left(B_{\ul{\nu}}\right)_{\ul{\nu}}$, with
\[
B_{\ul{\nu}} \;=\; \sum_{[\ul{\epsilon}] \in \mathrm{Mat}_{\ul{\Lambda}}(A;|\ul{\eta}|\times|\ul{\nu}|)}
\!\!\left(\prod^s_{i=1}c(\eta^i;R_i[\ul{\epsilon}])\right)\!\!\!\left(\prod^t_{j=1}c(\nu^j;C_j[\ul{\epsilon}])\right)
\]
where we define $A$ to be the $s \times t$ integer matrix whose $(i,j)$\tss{th} entry is $a^i_j$ (since $a^i_j=0$ is allowed, we note
that the empty multipartition $()$ may occur as an entry of $[\ul{\epsilon}]$). Further, suppose that we have $s = t$
and moreover for $i=1,\ldots,t$, $X_i$ occurs at the bottom of the filtration of $Y_i$. Then the module occurring at the bottom of this
filtration is $S^{\ul{\eta}}(X_1,\ldots,X_t)$.
\end{prop}

We remark that Proposition \ref{S_eta_Y_filt_S_nu_X:prop} is rather more general that is strictly required for proving the branching
rules below, but we feel that it is worth stating the result in this more general form, as it is of interest in its own right.

